\input{header.tex}
\input{newcommands.tex}


\title{Sitzpläne automatisch erstellen}
\author{Fortbildungsmaterial für Lehrkräfte}
\date{\today}
\pagestyle{plain}
\setlength{\parskip}{1.5em} % Abstand zwischen Absätzen (Handreichungen)
\setlength{\emergencystretch}{4em} % lange \texttt-Wörter umbrechen statt über den Rand


\begin{document}

\maketitle

Ein Sitzplan muss mehr können als Namen in Kästchen: Zwei Schüler sollen nicht nebeneinandersitzen, eine Schülerin braucht einen Platz vorn, ein Schüler arbeitet besser allein, und nach dem Halbjahr soll neu gemischt werden. Das Python-Skript \texttt{sitzplan.py} baut den Raum aus einer kurzen Beschreibung, verteilt die Namen per Zufall, bis alle Regeln erfüllt sind, und schreibt einen \hologo{LaTeX}-Plan in Lehrer- und Klassensicht. Ohne Namensliste entsteht ein leerer Plan zum Ausfüllen.

\begin{merke}
	\textbf{Datenschutz:} Schülernamen, Fotos mit Schülern und Noten gehören nicht in eine KI. Deshalb stehen die Namen in einer \textbf{eigenen CSV-Datei auf dem eigenen Rechner}. Eine KI darf beim Beschreiben des Raums helfen; die Namen setzt das Skript lokal ein. Die Beispieldateien enthalten nur Namen berühmter Mathematikerinnen, Mathematiker und Informatiker.
\end{merke}

\tableofcontents



%%%%  KAPITEL 1  %%%%
\section{Das Prinzip}

\begin{center}
\begin{tabular}{@{}>{\raggedright\arraybackslash}p{4.6cm}>{\raggedright\arraybackslash}p{3.6cm}>{\raggedright\arraybackslash}p{6.6cm}@{}}
	\toprule
	\textbf{Du änderst} & \textbf{Werkzeug} & \textbf{Ergebnis} \\
	\midrule
	\texttt{sitzplan\_\dots.toml} (Raum, Anordnung, Regeln) und eine CSV mit Namen & \texttt{sitzplan.py} & \texttt{sitzplan\_\dots.tex}, daraus mit \texttt{pdflatex} der Plan im Querformat \\
	\bottomrule
\end{tabular}
\end{center}

\begin{lstlisting}[style=latex]
python3 sitzplan.py sitzplan_10b.toml
pdflatex sitzplan_10b.tex
\end{lstlisting}

Benötigt wird Python ab Version~3.11, weitere Pakete sind nicht nötig. Das Skript meldet im Terminal, wie viele Namen und Plätze es gibt und ob Regeln verletzt werden mussten:

\begin{lstlisting}[style=latex]
geschrieben: sitzplan_10b.tex  (22 Namen, 24 Plätze, 0 Regelverstöße)
\end{lstlisting}



%%%%  KAPITEL 2  %%%%
\section{Die Namensliste}

Die Namen stehen in einer CSV-Datei mit den Spalten \texttt{Vorname;Nachname}, wie man sie aus der Schulverwaltung exportieren kann:

\begin{lstlisting}[style=latex]
Vorname;Nachname
Ada;Lovelace
Alan;Turing
Grace;Hopper
Konrad;Zuse
...
\end{lstlisting}

In den Regeln heißen die Personen dann \texttt{"Vorname Nachname"}, zum Beispiel \texttt{"Ada Lovelace"}.



%%%%  KAPITEL 3  %%%%
\section{Der Raum}


%%%%  UNTERKAPITEL 3.1  %%%%
\subsection{Allgemeine Angaben}

\begin{lstlisting}[style=latex]
titel = "Klasse 10b – Informatik"
stand = "November 2026"

[raum]
name        = "Raum 2.14"
ansichten   = ["lehrer", "schueler"]   # eine Seite je Ansicht
tafelbreite = 8
pult        = "links"                  # aus Klassensicht; "" = keins
fenster     = "links"
tuer        = "rechts"
\end{lstlisting}

\texttt{ansichten} legt fest, welche Seiten entstehen: \texttt{"lehrer"} mit der Tafel unten (für das Pult), \texttt{"schueler"} mit der Tafel oben (zum Aushängen).


%%%%  UNTERKAPITEL 3.2  %%%%
\subsection{Anordnungen}

Drei Anordnungen sind fertig eingebaut:

\begin{lstlisting}[style=latex]
[layout]                         # Reihen
typ               = "reihen"
reihen            = 4
tische_pro_reihe  = 3
plaetze_pro_tisch = 2

[layout]                         # Gruppentische
typ                = "gruppen"
gruppen_reihen     = 2
gruppen_pro_reihe  = 3
plaetze_pro_gruppe = 4           # 4 oder 6

[layout]                         # U-Form
typ    = "u_form"
links  = 4
hinten = 5
rechts = 4
\end{lstlisting}

Zusätzlich lassen sich einzelne Tische frei setzen; \texttt{x} und \texttt{y} sind in cm, $y = 0$ ist die Tafel:

\begin{lstlisting}[style=latex]
[[tisch]]
x = 20.5
y = 9
plaetze = 1
\end{lstlisting}



%%%%  KAPITEL 4  %%%%
\section{Namen verteilen und Regeln}

\begin{lstlisting}[style=latex]
[namen]
csv        = "klasse_beispiel.csv"   # bleibt auf deinem Rechner
verteilung = "zufall"                # "liste" = in CSV-Reihenfolge
seed       = 7                       # gleiche Zahl = gleicher Plan

[regeln]
nicht_neben  = [ ["Alan Turing", "Konrad Zuse"], ["Isaac Newton", "Gottfried Wilhelm Leibniz"] ]
neben        = [ ["Ada Lovelace", "Grace Hopper"] ]
vorne        = [ "Blaise Pascal" ]   # muss in die vorderste Reihe
vorne_reihen = 1
allein       = [ "Donald Knuth" ]    # Tisch für sich allein
[regeln.fest]
"Marie Curie" = 12                    # feste Platznummer
\end{lstlisting}

Das Skript mischt so lange, bis alle Regeln erfüllt sind, und meldet, wenn das nicht geht. Mit einem anderen \texttt{seed} entsteht eine neue Mischung, die wieder alle Regeln erfüllt; mit demselben \texttt{seed} entsteht immer derselbe Plan.

Das Ergebnis, links die Lehrersicht, rechts die Klassensicht:

\begin{ausgabe}
	\centering
	\fbox{\includegraphics[page=1, width=0.47\linewidth]{sitzplan_10b.pdf}}\hfill
	\fbox{\includegraphics[page=2, width=0.47\linewidth]{sitzplan_10b.pdf}}
\end{ausgabe}



%%%%  KAPITEL 5  %%%%
\section{Weitere Beispiele}

\texttt{sitzplan\_u\_form.toml} setzt einen Kurs mit 13 Personen in U-Form, in der Reihenfolge der CSV (\texttt{verteilung = "liste"}). \texttt{sitzplan\_gruppen\_leer.toml} hat keine CSV (\texttt{csv = ""}) und erzeugt einen leeren Plan mit Gruppentischen zum Ausfüllen von Hand:

\begin{ausgabe}
	\centering
	\fbox{\includegraphics[page=1, width=0.47\linewidth]{sitzplan_u_form.pdf}}\hfill
	\fbox{\includegraphics[page=1, width=0.47\linewidth]{sitzplan_gruppen_leer.pdf}}
\end{ausgabe}



%%%%  KAPITEL 6  %%%%
\section{Den Raum mit einer KI beschreiben}

Den Raum kann eine KI aus einem Foto beschreiben, sofern darauf \textbf{keine Personen} zu sehen sind. Die Namen bleiben dabei auf dem eigenen Rechner.

\begin{hinweis}[Beispiel-Prompt]
	\itshape Im Anhang ist ein Foto meines Klassenraums ohne Personen. Beschreibe den Raum als sitzplan.toml im Format der angehängten Vorlage: Tische als Layout oder einzelne [[tisch]]-Einträge (x, y in cm, y = 0 an der Tafel), Fenster, Tür und Lehrertisch. Keine Namen.
\end{hinweis}



%%%%  KAPITEL 7  %%%%
\section{Stolpersteine}

\begin{itemize}
	\item \textbf{Regel greift nicht?} Die Namen in den Regeln müssen genau wie in der CSV geschrieben sein, als \texttt{"Vorname Nachname"}.
	\item \textbf{\enquote{Regeln ließen sich nicht erfüllen}:} Die Regeln widersprechen sich oder es gibt zu wenige passende Plätze (etwa mehr Personen für \texttt{vorne} als Plätze in der vordersten Reihe). Der Plan wird trotzdem geschrieben und enthält einen roten Hinweis.
	\item \textbf{Mehr Namen als Plätze?} Weitere Tische mit \texttt{[[tisch]]} ergänzen oder die Anordnung vergrößern.
	\item \textbf{CSV mit Umlauten falsch?} Die Datei als UTF-8 speichern (in Tabellenkalkulationen beim Export wählbar).
\end{itemize}



%%%%  KAPITEL 8  %%%%
\section{Zusammenfassung}

\begin{itemize}
	\item Eine TOML-Datei beschreibt Raum, Anordnung und Regeln, eine CSV enthält die Namen.
	\item \texttt{python3 sitzplan.py plan.toml}, dann \texttt{pdflatex}: Plan in Lehrer- und Klassensicht.
	\item Regeln: \texttt{nicht\_neben}, \texttt{neben}, \texttt{vorne}, \texttt{allein}, \texttt{fest}; ein neuer \texttt{seed} ergibt eine neue Mischung.
	\item Namen gehören nie in eine KI; das Skript setzt sie lokal ein.
\end{itemize}

\end{document}
